% GENERATED FILE. Edit canonical lessons, then run npm run generate:books.
% canonical-source-hash: fnv1a64:1c20c753c60c392d
% canonical-lessons: TA-C97-definitely, TA-C97-thatshow, TA-C97-whynot

\chapter{Definitely, and Why Not}
\label{ch:ta-definitely}
\hlchaptermodality{97}

\begin{chapteropening}
\textbf{By the end of this chapter:} \emph{I can say definitely, that's how it is, and why not.}

Complete the last lesson of chapter 97: I can say definitely, that's how it is, and why not.
\end{chapteropening}

\section[\texorpdfstring{\ta{கண்டிப்பாக} (kaṇṭippāka)}{கண்டிப்பாக (kaṇṭippāka)}]{\ta{கண்டிப்பாக} (kaṇṭippāka) — definitely}

\label{lesson:TA-C97-definitely}



\begin{quote}
Before the new one: say the Tamil for that's right, then the Tamil for I don't know.
\end{quote}

\subsection*{You'll want to know: \ta{கண்டிப்பாக}}
\textbf{\ta{கண்டிப்பாக}} (\emph{kaṇṭippāka}) — \textquotedblleft{}definitely\textquotedblright{}, \textquotedblleft{}for sure\textquotedblright{}.

\textbf{\ta{கண்டிப்பாக} \ta{வருகிறேன்}} — \textquotedblleft{}I'll definitely come\textquotedblright{}.

\begin{cousinweb}
\textbf{\ta{கண்டிப்பாக}} (\emph{kaṇṭippāka}) is built on \emph{kaṇṭippu}, \textquotedblleft{}strictness, firmness\textquotedblright{}, which is usually traced to Sanskrit \emph{khaṇḍ-}, \textquotedblleft{}to break, to cut, to settle\textquotedblright{} (CDIAL 3797, \emph{khaṇḍita}, \textquotedblleft{}broken\textquotedblright{}). \emph{-āka} makes it \textquotedblleft{}definitely\textquotedblright{}.
\end{cousinweb}

\begin{grammarlens}[title={What you've built}]
A firm yes.
\end{grammarlens}

\subsection*{Your turn}
\begin{itemize}
\raggedright
  \item \emph{Say it:} \emph{kaṇṭippāka}
  \item \emph{Say it:} \emph{kaṇṭippāka}, once more
  \item \emph{Say it:} say \emph{kaṇṭippāka}
  \item \emph{Recall:} say the Tamil for that's right, then the Tamil for I don't know, then say \emph{kaṇṭippāka} again
\end{itemize}

\begin{tcolorbox}[breakable,colback=teal!4,colframe=teal!35!black,title={Before you move on}]
What does \ta{கண்டிப்பாக} mean? (\textquotedblleft{}Definitely\textquotedblright{}.) Say it once more.
\end{tcolorbox}
\section[\texorpdfstring{\ta{அப்படித்தான்} (appaṭittāṉ)}{அப்படித்தான் (appaṭittāṉ)}]{\ta{அப்படித்தான்} (appaṭittāṉ) — that's how it is}

\label{lesson:TA-C97-thatshow}



\begin{quote}
Before the new one: say the Tamil for I don't know, then the Tamil for definitely.
\end{quote}

\subsection*{You'll want to know: \ta{அப்படித்தான்}}
\textbf{\ta{அப்படித்தான்}} (\emph{appaṭittāṉ}) — \textquotedblleft{}that's how it is\textquotedblright{}, \textquotedblleft{}quite so\textquotedblright{}.

\textbf{\ta{அப்படி}}, \textquotedblleft{}like that\textquotedblright{}, with \textbf{-\ta{தான்}}, the same \textquotedblleft{}that and no other\textquotedblright{} as \textbf{\ta{சரிதான்}}.

\begin{cousinweb}
\textbf{\ta{அப்படித்தான்}} (\emph{appaṭittāṉ}) is native: \emph{appaṭi}, \textquotedblleft{}thus\textquotedblright{} (DEDR 1), and \emph{tāṉ}, \textquotedblleft{}just, indeed\textquotedblright{} (DEDR 3196).
\end{cousinweb}

\begin{grammarlens}[title={What you've built}]
A shrug that agrees.
\end{grammarlens}

\subsection*{Your turn}
\begin{itemize}
\raggedright
  \item \emph{Say it:} \emph{appaṭittāṉ}
  \item \emph{Say it:} \emph{appaṭittāṉ}, once more
  \item \emph{Say it:} say \emph{appaṭittāṉ}
  \item \emph{Recall:} say the Tamil for I don't know, then the Tamil for definitely, then say \emph{appaṭittāṉ} again
\end{itemize}

\begin{tcolorbox}[breakable,colback=teal!4,colframe=teal!35!black,title={Before you move on}]
What does \ta{அப்படித்தான்} mean? (\textquotedblleft{}That's how it is\textquotedblright{}.) Say it once more.
\end{tcolorbox}
\section[\texorpdfstring{\ta{ஏன்} \ta{இல்லை}? (ēṉ illai?)}{ஏன் இல்லை? (ēṉ illai?)}]{\ta{ஏன்} \ta{இல்லை}? (ēṉ illai?) — why not?}

\label{lesson:TA-C97-whynot}



\begin{quote}
Before the new one: say the Tamil for definitely, then the Tamil for that's how it is.
\end{quote}

\subsection*{You'll want to know: \ta{ஏன்} \ta{இல்லை}?}
\textbf{\ta{ஏன்} \ta{இல்லை}?} (\emph{ēṉ illai?}) — \textquotedblleft{}why not?\textquotedblright{}.

\textbf{\ta{ஏன்}} is \textquotedblleft{}why\textquotedblright{}; with \textbf{\ta{இல்லை}}, \textquotedblleft{}why not?\textquotedblright{} — as a question, or as a cheerful yes.

\begin{cousinweb}
\textbf{\ta{ஏன்} \ta{இல்லை}?} (\emph{ēṉ illai?}) is native: \emph{ēṉ}, \textquotedblleft{}why\textquotedblright{} (DEDR 5151), and \emph{illai}, \textquotedblleft{}not\textquotedblright{} (DEDR 2559).
\end{cousinweb}

\begin{grammarlens}[title={What you've built}]
Definitely, that's how it is, why not.
\end{grammarlens}

\subsection*{Your turn}
\begin{itemize}
\raggedright
  \item \emph{Say it:} \emph{ēṉ illai?}
  \item \emph{Say it:} \emph{ēṉ illai?}, once more
  \item \emph{Say it:} say \emph{ēṉ illai?}
  \item \emph{Recall:} say the Tamil for definitely, then the Tamil for that's how it is, then say \emph{ēṉ illai?} again
\end{itemize}

\begin{tcolorbox}[breakable,colback=teal!4,colframe=teal!35!black,title={Before you move on}]
What does \ta{ஏன்} \ta{இல்லை}? mean? (\textquotedblleft{}Why not?\textquotedblright{}.) Say it once more.
\end{tcolorbox}
